\documentclass[sigplan,review,anonymous]{acmart}

\usepackage{minted}
\usepackage{graphicx}
\usepackage{subcaption}
\usepackage{cleveref}
\usepackage[inline]{enumitem}

\title{Sulfur: Substitution Generation using a Logical Framework}

\author{
  Mathis Bouverot-Dupuis
  \and
  Théo Winterhalter
}

\newcommand{\coq}[1]{\mintinline{coq}{#1}}

\setminted{fontsize=\footnotesize}

\begin{document}

\maketitle

\section{Introduction}

When formalizing the meta-theory of programming languages or type systems
in proof assistants
such as Rocq~\cite{rocq}, a key design decision is the choice of
variable representation.
%Many options are available:
%fully named representations (used e.g. in Nominal
% Isabelle~\cite{nominal-isabelle})
%which stay close to informal practice, higher-order abstract syntax
%(HOAS)~\citet{hoas-pfenning}
%(used e.g. in Twelf~\cite{twelf}) and variats such as parametric
%higher-order abstract syntax (PHOAS)~\cite{phoas}, nameless representations
%such as de Bruijn indices~\cite{de-bruijn-indices} and well-scoped
% variants~\cite{},
%or the locally nameless representation
%% many options:
%% - named representations (nominal isabelle)
%% - HOAS (twelf? what else?)
%% - locally nameless
%% - de Bruijn indices
De Bruijn indices~\cite{de-bruijn-indices} are a popular option
because they make alpha-equivalence
coincide with syntactic equality, and because sigma
calculus~\cite{sigma-calculus, sigma-calculus-completeness} provides a
complete equational theory to reason about substitution.
However, de Bruijn indices also introduce significant
overhead in the form of lift and renaming operations, which require many
technical lemmas and can make proofs tedious.

To address this, tools such as
Autosubst~\cite{autosubst-1,autosubst-2} automate
repetitive aspects of
dealing with de Bruijn indices and substitution using boilerplate generation.
However, Autosubst can have significant performance issues when used in large
developments.
In this talk, I present Sulfur, a Rocq plugin that similarly
automates reasoning
about substitutions, and attemps to solve some of the performance
issues of Autosubst.
The user interface of Sulfur is intentionally very close to that
of Autosubst,
however the implementation diverges significantly.

\section{Using Sulfur}

Given a user-defined language signature (\Cref{fig:sulfur-example-a}), Sulfur
automatically generates Rocq code (\Cref{fig:sulfur-example-b})
implementing an inductive type representing terms
with variables encoded as de Bruijn indices, and a parallel
substitution function.

\begin{figure}[H]
  \centering
  \begin{subfigure}{0.4\textwidth}
\begin{minted}{coq}
term : Type
app : term -> term -> term
lam : (bind term in term) -> term
\end{minted}
    \subcaption{User-specified signature.}
    \label{fig:sulfur-example-a}
  \end{subfigure}

  \begin{subfigure}{0.4\textwidth}
  \begin{minted}{coq}
    Inductive term :=
    | var (i : nat)
    | app (t u : term)
    | lam (t : term).

    Definition substitute : (nat -> term) -> term -> term.
    (** substitute s t is abbreviated as t[s]. *)
    \end{minted}
    \subcaption{Code generated by Sulfur.}
    \label{fig:sulfur-example-b}
  \end{subfigure}

  \caption{Untyped lambda calculus, using Sulfur}
  \label{fig:sulfur-example}
\end{figure}

Most importantly, Sulfur automates reasoning about substitution:
it provides a tactic \coq{asimpl} which simplifies terms and substitutions
according to the rules of sigma calculus. For instance consider
the following technical lemma, which is needed when proving
standard properties of lambda calculus:

\begin{minted}{coq}
Lemma technical_lemma (t1 t2 : term) (s : nat -> term) :
  t1[0 . (s >> shift)][t2[s] . id] = t1[t2 . id][s].
\end{minted}
where \coq{id} is the identity substitution,
\coq{shift} is the substitution which adds 1 to every de
Bruijn index, \coq{t . s} is a substitution which maps index \coq{0} to \coq{t}
and index \coq{i+1} to \coq{s i}, and \coq{s1 >> s2} is the
composition of \coq{s1} followed by \coq{s2}.

Proving this equality by hand requires significantly more effort
than one might
expect, requiring many auxiliary lemmas. However,
\coq{asimpl} simplifies both sides of the equation to
\coq{t1[t2[s] . s]}, from which the proof can be trivially finished.

\section{Key ideas}

While Sulfur provides the same user interface as Autosubst,
its implementation diverges significantly. In particular, Autosubst's
\coq{asimpl} tactic relies on Rocq's \textit{setoid rewrite}
facilities, and has significant
performance issues when used in large developments.
Sulfur aims to improve performance by using a logical framework approach:
we define a generic notion of syntax with explicit substitutions
within Rocq, and implement \coq{asimpl} as a plain Rocq function over
this generic syntax.

\paragraph{Generic syntax}

% 1. concrete syntax (which is what the user sees): see code above
% 2. logical framework: generic notion of syntax w/ explicit substitutions

Inspired by sigma calclus, we define a notion of syntax with explicit
substitutions, which is moreover \textit{generic}, i.e. parameterized over
a signature. We give a simplified version of this syntax:
\begin{minted}{coq}
Inductive gterm sig :=
| gvar (i : nat)
| gctor (c : ctor sig) (args : list gterm)
| gsubstitute (s : gsubst) (t : gterm)
with gsubst sig :=
| gid
| gshift
| gcons (t : gterm) (s : gsubst)
| gcomp (s1 s2 : gsubst).
\end{minted}

Generic terms \coq{gterm} and generic substitutions \coq{gsubst} are
parameterized
over a signature \coq{sig}, which encodes the set of constructors
\coq{ctor sig}
along with information about the arity and binding structure of each
constructor. Notably, substitutions are not arbitrary functions
\coq{nat -> gterm},
but instead are built using a set of constructors \coq{gid},
\coq{gshift}, \coq{gcons}, and \coq{gcomp}, which correspond to
the functions \coq{id}, \coq{shift}, \coq{_ . _}, and \coq{_ >> _}.
Moreover, a substitution can be explicitly applied to a term using
\coq{gsubstitute}.

Generic syntax contains enough information to be able
to implement a simplification function directly in Rocq:
\begin{minted}{coq}
Definition simplify : forall sig, gterm sig -> gterm sig.
\end{minted}

\paragraph{Reification and denotation}

Generic syntax is quite far from what we picture as the untyped
lambda calculus, and we certainly
do not want users to work with \coq{gterm}. Thus, we still generate
syntax specialized to the user's signature, as in~\Cref{fig:sulfur-example-b}.

The mapping between user syntax and generic syntax, i.e. between \coq{term}
and \coq{gterm}, is fairly straightforward.
A \textit{denotation} function
\coq{denote} from \coq{gterm sig} to \coq{term} (where \coq{sig} is a signature
encoding the untyped lambda calculus) can be implemented directly within Rocq,
by simple structural recursion over its argument. The other direction,
which we call \textit{reification}, requires us to step outside of Rocq
and leverage Rocq's support for meta-programming. Reification and
denotation are required to
obey the invariant that for any \coq{t : term}, if reifying \coq{t} yields
\coq{t' : gterm sig}, then \coq{denote t'} is convertible to \coq{t}.

\paragraph{Soundness}

In order to simplify terms in user-syntax, we need to establish the
soundness of \coq{simplify}:

\begin{minted}{coq}
Theorem soundness :
  forall t : gterm sig, denote t = denote (simplify t).
\end{minted}

% 5. which is proved sound
% w.r.t. concrete syntax. This is done in two steps:
% - I define a reduction system on generic syntax, which is proved sound
%   in the sense that [red t t' -> denote t = denote t']
%\begin{minted}{coq}
%
%\end{minted}
% - I write a simplification function which chooses an order in which
%   to apply the reduction rules: [simpl : term -> term], with a spec
%   [red t (simpl t)].

\paragraph{Putting it all together}

Using all these ingredients, we can implement \coq{asimpl} as
follows. To simplify a term \coq{t : term} appearing in a goal:
\begin{itemize}
  \item Reify \coq{t} into \coq{t' : gterm sig}.
  \item Because \coq{t} is convertible to \coq{denote t'}, which itself is equal
    to \coq{denote (simplify t')}, we can replace all occurences of
    \coq{t} with \coq{denote (simplify t')} in the goal.
  \item Use Rocq's evaluation mechanisms to reduce \coq{denote (simplify t')}.
\end{itemize}

In terms of efficiency the key benefit of this approach is that,
contrary to setoid rewriting,
\coq{asimpl} does not build a proof of equality (between the original
term \coq{t} and its simplified version) on the fly: we instead prove
a soundness theorem once and for all.

%We will discuss the tradeoffs of the logical framework approach,
%in particular regarding implementation effort and performance of `asimpl`.
%

\section{Related work}

% on substitution generation: many libraries, both in mainstream
% programming languages (e.g AlphaCaml, FreshML, Unbound, Rebound) and in
% proof assistants (e.g. Autosubst, LNGen, Tealeaves, Fiore's paper).
% Most libraries focus either on performance of the generated substitution
% functions, or on generating many useful lemmas.
% With Sulfur, we instead focus on implementing
% an efficient simplification tactic.

% on reification tactics: this is a classic idea, used to
% implement many tactics. A peculiarity of Sulfur is that
% the target language is very complex, encoding intricate
% invariant using dependent types.

% on

\section{Future work}

\paragraph{More complex signatures.} Central future work is to extend
Sulfur to accomodate
more complex signatures, scaling up to the full generality of
Autosubst. For instance, signatures with multiple sorts of terms
(e.g. System F) and signatures including functors (e.g. using lists to represent
n-ary applications) are not supported yet. We believe that extending Sulfur
with these features is not simply a matter of engineering work, and will require
novel insight.

\paragraph{Proving completeness.} The following completeness theorem holds in
simpler variants of sigma calculus~\cite{sigma-calculus-completeness}:
\begin{minted}{coq}
Theorem completeness (t t' : gterm sig) :
  denote t = denote t' -> simplify t = simplify t'
\end{minted}
Intuitively this states that reification followed by simplification
is enough to decide equality of concrete terms.
Unfortunately, this completeness theorem does not hold on our version
of generic syntax, however we conjecture that some weaker version of
completeness still holds, and believe that proving such a result
is an interesting direction for future work.

\bibliographystyle{ACM-Reference-Format}
\bibliography{biblio}

\end{document}